\subsection{Complete windows along finite histories}\label{sec:complete-windows}
Fix an actual nonempty column prefix $U$ in a column $q$, let $p$ be
its last parent, and let $s$ be the next parent, or the parentless
sentinel if there is no next entry. A \emph{qualified cut} consists of
a root $r>s$ on the single residual-parent chain of $p$ for the entire
word $\mathcal H(U)$. Let $X$ be the complete same-level suffix beginning
at this cut. It includes all later columns and their private heights.

Temporarily truncate to $U$ and append a formal bare entry with parent
$r$. Apply one native copy window while leaving its rightmost recursive
hole unfilled. This defines a one-hole context $\Psi$. It retains $U$,
appends the complete root residue, shifts addresses by the native
threshold, and fills the hole by its argument. If $r$ contains $q$,
property G makes the whole threshold word strictly lower than the
containing head. Hence the scan stops before that head is consumed,
and the ancestor branch of the native formula still defines the same
one-hole context. The formal context is only used for comparison; it
is not asserted to be an actual node.

The hole follows a nonempty finite constructor. Both ordinary and
ancestor contexts strictly preserve order and recursive prefixes.
Define $T_0=\varnothing$ and $T_{j+1}=\Psi(T_j)$.

\begin{lemma}[Finite context criterion]\label{lem:context}
The $T_j$ form an increasing chain of recursive prefixes, and for each
finite complete suffix $X$,
\begin{equation}\label{eq:3-3}
       X<\Psi(X)\quad\Longleftrightarrow\quad
       (\exists j\ge1)\ X\le T_j.                    
\end{equation}
\end{lemma}
\begin{proof}
The empty suffix is a prefix of $\Psi(\varnothing)$; prefix preservation
proves the first assertion. If $X\le T_j$, compare $X$ with $\Psi(X)$
at the first place where $X$ leaves the common iterated framework.
Leaving earlier gives $X<\Psi(X)$, while staying in the framework
reduces the comparison through the hole. After at most $j$ reductions
one reaches the empty suffix, where strictness holds. Conversely,
suppose all $T_j<X$. The first difference between $X$ and each $T_j$
must lie after the framework shared with the preceding iterate; otherwise
strict order preservation would contradict $X<\Psi(X)$ at that same
difference. Each reduction passes a nonempty constructor. Eventually
the common framework would pass the last field of finite $X$, which
is impossible. This is induction on the finite first-difference path,
not on BTBMS descent.
\end{proof}

\begin{lemma}[Native window transfer]\label{lem:window-transfer}
For corresponding cuts in copies of the same mother step, let $\sigma$
be their actual address transport and suppose its threshold $a$ satisfies
$r\le a$ for the fixed root $r$. With the native remaining-copy indices,
\begin{equation}\label{eq:3-4}
 \sigma\Psi_i(X_i)\pref\Psi_j(X_j),\qquad
 X_j\le\sigma X_i.                                   
\end{equation}
\end{lemma}
\begin{proof}
Pull the finite comparison word back separately in each copy. Before an
ancestor is met, closed heads are compared by the same order-preserving
map. At the ancestor, G stops the scan before its changing tail. After
the first transport, the root's private addresses are internal, and
the second shift adds to the first. True external parents and the top
parents of the root residue remain fixed. Thus field-by-field transport
commutes with filling the complete hole; the possible loss is only a
right truncation. This proves (\ref{eq:3-4}).

If the source mother root itself is open, copy $j$ has only $n-j$
mother iterations left. Using $n$ would fill the wrong hole. There is
also a separate owner case: if the mother tail is $t$, the root is $b$
and $d=t-b$, then the $j$th transported tail-prefix package has owner
$b+(j+1)d$, not $b+jd$. A query to its own containing owner uses the
closed containing package and commutes with transport. A query to an
earlier boundary either leaves a head in the installed prefix, in which
case its first residual parent is already strictly above the next parent
$s$, or exhausts that prefix and returns to the mother's residue.
In the latter case apply the residual identity of Lemma~\ref{lem:word}
and the earlier-boundary instance of (\ref{eq:3-4}). The boundary index decreases.
This treats owner queries without applying (\ref{eq:3-4}) at an ineligible
threshold below its fixed root.
\end{proof}

\begin{proposition}[Actual complete-window bound]\label{prop:sce}
Every qualified cut of every $M\in\mathscr U$ satisfies $X<\Psi(X)$.
Every standard completed candidate satisfies the corresponding statement
for roots internal to that candidate.
\end{proposition}
\begin{proof}
Induct on the length of the given finite native history. In an initial
tower, a cut lies at the end of a singleton column. Its suffix is
already a prefix of the first nonempty window, so
Lemma~\ref{lem:context} applies. For the induction step classify the
position of the cut in the explicit native construction.
\begin{enumerate}
\item At a retained prefix $U$, closed earlier roots and the formal
window are unchanged. An open containing root is protected by G. The
suffix only decreases, so its old finite-iterate bound remains. If the
last entry of $U$ changes internally, the relevant level has the single
retained column $U$, which is a prefix of $\Psi(\varnothing)$.
\item At the first new boundary, EA puts the old next parent and the
new residual parent consecutively on the same chain. The only new
qualified root between them is the mother root. Its new suffix is the
actual $T_n$, giving the required finite bound directly.
\item A cut in the root residue uses the same normal form of the
complete scanned prefix. Exhausting an installed prefix composes the
same one-hole context. If its final retained entry changes, use the
single-column case of the first item.
\item A cut in the bad part, a root-private package or a later boundary
is governed by Lemma~\ref{lem:window-transfer}. A query leaving a head
in the installed prefix is decided by the strict first-parent comparison.
An exhausting query uses the old residual window and the correct
remaining-copy index. Moving to an earlier boundary decreases that
finite index.
\item At the final zero hole, the entire deepest last column disappears.
A removed private cut has no surviving obligation. A retained cut
whose $U$ changes is the single-column case; otherwise its containing
head is protected. Special repetitions preserve same-parent adjacency;
later copies use the preceding complete copy and the same window fill.
A new same-parent seam has no new root satisfying $r>s$.
\end{enumerate}
These cases cover the retained prefix, new seam, residue, bad part,
private copies and deleted tail. They use the induction hypothesis only
on the actual mother node, not on an arbitrary formal substitution.

For Cut, restrict attention first to an internal root. If Cut changes
the last retained entry of $U$, the output has the single-column form
already considered. Otherwise all earlier closed fields are unchanged,
and G protects every containing head before the cut. The same formal
context has legal finite iterates at the destination owner: its retained
framework is legal, and copying only moves internal fields. Choose a
strict old finite-iterate upper bound. Minimality of Cut bounds the new
suffix by that legal iterate, and Lemma~\ref{lem:context} applies again.
Recursive prefixes, true subheights and whole address transport preserve
this argument. Birth-record induction therefore gives the assertion for
all standard completed candidates. No arbitrary external port is added
by this internal version.
\end{proof}

\subsection{Source ports and strict carriers}\label{sec:strict-proof}
There are two transport issues in applying Proposition~\ref{prop:sce}.
First, equal numerical addresses in unrelated private environments are
not common ports. Second, a Cut-created bare edge has to be pulled back
to an old chain; it cannot be declared an old edge.

Suppose a whole package at owner $q$ is moved to a compatible port $a$.
For a retained non-bare internal entry, transport an arbitrary finite
threshold word of smaller depth by its finite comparison pattern.
Closed fields before Cut are either true external parents below $a$ or
internal fields at least $q$; transport is strictly increasing on this
allowed set. The standard heights in a common external column are
fixed by the same comparison map. At a containing ancestor G again
prevents the scan from reading the changed tail. If the new chain goes
through the new owner, the old chain goes through $q$, and the common-port
condition supplies the rest of that very same word's chain to $a$.
If it uses the new final bare edge, (\ref{eq:2-1}) identifies its word with the
normal form of the cut prefix. Old EA then supplies the missing old
segment. If the bare entry is consumed, there is no new residual edge
to justify. This proves that every compatible common external port of
the moved entry was already a compatible port of its actual source.

Tracing a non-bare head $B$ in an original completed column $r$ through
its actual birth record now gives a source occurrence $x$, of height
$X$, such that
\begin{equation}\label{eq:3-5}
 B=C_r(X),\quad d(B)=d(X)=g,\quad
 r\in\{\operatorname{owner}(x)\}\cup
       \Pi_g(\operatorname{owner}(x)).                 
\end{equation}
This traces one actual column, including through private copies; it
does not infer provenance from a printed parent number.

New support in the original gap immediately before an active height
$H$ has depth below $d(H)$. Indeed a request from within $H$ has that
depth bound by G; a request from an earlier closed package terminates
before its own containing head and never enters this gap; a request
from a package newly inserted in the gap has still lower depth. Induct
on creation order. Therefore such support lies below every positive
result $P$ with $d(P)=d(H)$ and cannot be a strict barrier.

\begin{proposition}[Strict-barrier theorem]\label{prop:strict-barrier}
In an actual corresponding backward query for $H\to P$, every original
strict carrier $P<B<H$ contains $P$ as a recursive prefix. The same
conclusion holds after the permitted preparation, with the four carrier
origins specified below.
\end{proposition}
\begin{proof}
Consider first a positive general change with nonempty deepest tail
prefix $U$. If the mother root $b$ is internal to $H$, the interval
$P<B<H$ forces $B$ to retain the entire framework through $U$, with
next parent below $b$. EA therefore leads to the same corresponding
internal root: a chain cannot leave the package and later return to
an internal column, and G protects any containing head on the route.
The two complete contexts agree. The suffix of $P$ is $T_n$; the suffix
$Y$ of $B$ has a finite bound $Y\le T_k$ by Proposition~\ref{prop:sce}.
Thus $T_n\le Y\le T_k$. Since the iterates are recursive prefixes,
prefix interval convexity gives $T_n\pref Y$, hence $P\spref B$.

If $b$ is external, let $r<q$ be the compatible actual query column.
Candidate dichotomy and minimality give $C_r(P)=P$. The owner-boundary
cases in that dichotomy exclude $b=q$ and $r<b<q$, so $b\le r<q$.
Equation (\ref{eq:3-5}), followed back to the first strict source, gives
\[
                    P<X\le B=C_r(X)<H.
\]
This $x$ is not the active entry or an entry inside $H$. The common
framework up to the deleted last bare entry identifies a cut in $X$
with the same complete $U$ and next parent below $b$. For internal
pieces, comparison is transported by the same order map. All true
external fields of that framework are below $r$ and hence name common
actual columns. If the chain leaves the source owner, (\ref{eq:3-5}) puts both
threshold chains through that same $r$; their tails then agree down to
$b$. Proposition~\ref{prop:sce} gives $Y<\Psi_x(Y)$ for the complete
source suffix $Y$.

Perform only the actual zero deletion at this source cut, and call its
remaining complete suffix $J$. The common framework survives: otherwise
it would contain the global deepest bare entry and have fewer entries
than $H$, contrary to $P<X<H$ at this cut. Put $m=\nu([U])$.
Then $J\pref Y$, and the removed suffix has size either zero or $m+1$.
The fixed part of $\Psi_x$ contains all of $U$ and one additional entry
leading to the hole, hence has size at least $m+1$. For an ancestor
root that extra entry survives by G. Filling the full window therefore
gives, under the actual source-target field correspondence,
\begin{equation}\label{eq:3-6}
               Y\le\widehat{\Psi_x(J)}\pref T_1\pref T_n.
\end{equation}
Here both internal shifts add and true external parents remain fixed;
one must retain all same-level columns when making this comparison.
But $P<X$ gives $T_n\le Y$. All comparisons in (\ref{eq:3-6}) must be equalities.
If $J=Y$, the source strict window bound contradicts equality. If
$J\ne Y$, equality leaves no unaccounted field outside the rightmost
containing chain and gives $X=P$, also a contradiction. Thus an external
root has no strict carrier in this interval.

If $U$ is empty, use the complete-column window: the root is an earlier
nonempty direct private column, its entire residue survives the empty
scan, and the preceding internal-root argument uses that full column
as the nonempty context. Zero changes are recursive truncations, and
special changes have repeated identical same-parent constructors; in
both cases the first-difference prefix comparison gives the assertion
directly. Later strict carriers inherit the first one's prefix by
recursive-prefix interval convexity. The four permitted preparation
origins are checked in the following paragraphs and
Lemma~\ref{lem:empty-prefix}.
\end{proof}
